\scanpage{100}
\uline{Dém} C'est géom., on peut se placer sur corps alg. clos. Nécessité : Supposons que ce soit Azumaya, \deleted{On peut d'} \inserted{prouvons non dégénérescence.}

Si $A_1,A_{\unsure{1|2}}$ sont unipotentes, on a prouvé le résultat, et aussi si $A_1,A_2$ sont toutes les deux étales. Il reste le cas où l'une est étale, l'autre unipotente, mais il n'y a qu'une seule position relative \inserted{possible}. Dans ce cas on choisit $v$ nilpotent, $u$ un projecteur
\[
\begin{aligned}
&\deleted{b=1,\ c=0},\quad b'=0,\ c'=0\\
&Q=t^2,\quad\Delta=-Q=-t^2
\end{aligned}
\]
: vérifier que $t\ne\unsure{0|10}$. On choisit
\[
\deleted{\begin{pmatrix}0&1\\0&0\end{pmatrix}=v,\quad u=\begin{pmatrix}1&0\\0&0\end{pmatrix}\quad uv=\begin{pmatrix}0&1\\0&0\end{pmatrix}}
\]
\[
\begin{aligned}
v&=\begin{pmatrix}0&1\\0&0\end{pmatrix}&u&=\begin{pmatrix}\alpha&\beta\\\gamma&\delta\end{pmatrix}\quad\text{``général'',}\\
uv&=\begin{pmatrix}0&\alpha\\0&\gamma\end{pmatrix}&\operatorname{Tr}uv&=\gamma\\
&&\delta(\unsure{u|x})&=(\alpha+\delta)^2-4(\alpha\delta-\beta\gamma)
\end{aligned}
\]
On peut choisir $\alpha\beta\gamma\delta$ tels que $\gamma\delta(\alpha,\beta,\gamma,\delta)\ne0$

Suffisance \deleted{Il suffit de} \ldots.
% archive sheet number: 99
